% !TeX root = ../main.tex

% ==================================================
% A.5. nonlocal과 클로저의 상태 변경
% ==================================================

\subsection{\zsh{nonlocal}과 클로저의 상태 변경}

Nonlocal Binding · Shared Closure Cells · Rebinding · Mutation · State Management

\subsubsection{이전 학습과의 연결}

\zsh{A.4}에서는 클로저(Closure)가 단순히 외부 변수의 값을 복사하는 구조가 아니라, 

렉시컬 외부 함수의 바인딩에 접근할 수 있는 환경을 유지하는 구조라는 사실을 학습했다.

특히 다음 객체 관계를 확인했다.

\begin{dhkbranch}[in=8]
node main :: 함수 객체
branchline main -> c :: __code__
handoff c -> c1 :: co_freevars
stop c1
handoff main -> d :: __closure__
arrow d
node d :: Cell 객체
arrow d
node d :: 참조하는 객체
\end{dhkbranch}

그리고 하나의 외부 함수 호출에서 생성된 여러 내부 함수가 동일한 자유 변수를 참조한다면, 같은 셀 객체를 공유할 수 있다는 사실도 확인했다.

그런데 아직 중요한 문제가 남아 있다.

외부 함수의 바인딩을 읽는 것과 그 바인딩을 변경하는 것은 어떻게 다를까?

다음 코드를 살펴보자.

\begin{dhkcode}
def outer():
    x = 10

    def get():
        return x

    def set_value(value):
        nonlocal x
        x = value

    return get, set_value


get, set_value = outer()

print(get())

set_value(100)

print(get())
\end{dhkcode}

예상 출력:

\begin{dhkzsh}
10
100
\end{dhkzsh}

여기서 \zsh{get}과 \zsh{set_value}는 서로 다른 함수 객체다.

그런데 \zsh{set_value()}를 호출하면 \zsh{get()}의 반환값까지 변경된다.

왜 이런 결과가 발생할까?

두 함수가 같은 셀 객체를 공유하고, \zsh{set_value()}의 \zsh{nonlocal x} 선언이 그 셀에 연결된 바인딩을 변경하기 때문이다.

이번 학습에서는 이 과정을 Python의 객체 모델에 따라 분석할 것이다.

\subsubsection{이번 학습의 핵심 질문}

이번 학습에서 해결할 질문은 다음과 같다.

\begin{dhksmry}
1. nonlocal은 정확히 어떤 이름을 대상으로 하는가?

2. nonlocal 선언은 셀 객체를 생성하거나 값을 변경하는가?

3. 여러 클로저가 동일한 셀을 공유하면 어떤 일이 발생하는가?

4. 셀 객체의 동일성과 셀이 참조하는 객체의 동일성은 어떻게 다른가?

5. 이름 재바인딩과 가변 객체의 상태 변경은 어떻게 구별하는가?

6. nonlocal 없이도 클로저의 상태를 변경할 수 있는가?

7. nonlocal은 왜 특정 상황에서 SyntaxError 또는
   UnboundLocalError와 연결되는가?

8. 클로저를 이용한 상태 관리는 클래스 인스턴스의 상태 관리와 어떻게 다른가?
\end{dhksmry}

이번 학습 전체를 관통하는 핵심 원리는 다음과 같다.

\begin{dhksmry}[centered]
nonlocal은 외부 객체를 변경하는 명령이 아니라, 현재 함수 블록의 

이름 바인딩 대상을 렉시컬 외부 함수 스코프로 지정하는 선언이다.
\end{dhksmry}

이 차이를 정확하게 이해해야 한다.

\subsubsection{\zsh{nonlocal}의 정확한 정의}

\zsh{nonlocal}은 현재 함수 블록에서 사용하는 이름이 현재 함수의 지역 이름이 아니라, 렉시컬 구조상 바깥에 있는 함수 스코프의 기존 바인딩을 사용하도록 지정하는 선언문이다.

가장 기본적인 형태는 다음과 같다.

\begin{dhkcode}
def outer():
    x = 10

    def inner():
        nonlocal x
        x = 20

    inner()

    return x


print(outer())
\end{dhkcode}

예상 출력 결과:

\begin{dhkzsh}
20
\end{dhkzsh}

여기서 \zsh{nonlocal x}가 수행하는 역할을 정확하게 구분해 보자.

\begin{dhknote}
\begin{dhksmry}[centered]
nonlocal x 자체가 수행하지 않는 일
\end{dhksmry}
\end{dhknote}

\zsh{nonlocal x}는 다음 작업을 직접 수행하지 않는다.

\begin{dhksmry}
    - 새로운 정수 객체를 생성하지 않는다.

    - 새로운 셀 객체를 매번 생성하지 않는다.

    - 외부 함수의 x를 즉시 변경하지 않는다.

    - 외부 함수의 지역 변수를 현재 함수로 복사하지 않는다.
\end{dhksmry}

대신 현재 함수에서 \zsh{x}라는 이름을 어떻게 해석하고 바인딩할지 지정한다.

실제 바인딩 변경은 다음 문장이 실행될 때 발생한다.

\begin{dhkcode}
x = 20
\end{dhkcode}

이때 \zsh{x}가 현재 함수의 지역 이름이 아니라 외부 함수의 바인딩을 대상으로 한다는 점이 중요하다.

\subsubsection{\zsh{nonlocal}이 없는 경우와 있는 경우}

먼저 두 코드를 비교해 보자.

\begin{dhknote}
\begin{dhksmry}[centered]
사례 1. nonlocal이 없는 경우
\end{dhksmry}
\end{dhknote}

\begin{dhkcode}
def outer():
    x = 10

    def inner():
        x = 20
        return x

    result = inner()

    print("inner:", result)
    print("outer:", x)


outer()
\end{dhkcode}

예상 출력:

\begin{dhkzsh}
inner: 20
outer: 10
\end{dhkzsh}

왜 외부 함수의 \zsh{x}는 변경되지 않았을까?

\zsh{A.3}에서 학습한 이름의 정적 판정 규칙 때문이다.

\zsh{inner} 내부에는 다음 대입문이 있다.

\begin{dhkcode}
x = 20
\end{dhkcode}

별도의 \zsh{global} 또는 \zsh{nonlocal} 선언이 없으므로 \zsh{x}는 \zsh{inner}의 지역 이름으로 판정된다.

따라서 두 함수에는 서로 다른 바인딩이 존재한다.

\begin{dhkdir}[in=6]
outer() 실행 환경
> x -> int(10)
\end{dhkdir}

\begin{dhkdir}[in=6]
inner() 실행 환경
> x -> int(20)
\end{dhkdir}

두 이름의 철자가 같더라도 동일한 바인딩을 의미하지 않는다.

따라서 \zsh{inner()}가 종료된 뒤에도 \zsh{outer}의 \zsh{x}는 10을 참조한다.

\begin{dhknote}
\begin{dhksmry}[centered]
사례 2. nonlocal이 있는 경우
\end{dhksmry}
\end{dhknote}

\begin{dhkcode}
def outer():
    x = 10

    def inner():
        nonlocal x
        x = 20
        return x

    result = inner()

    print("inner:", result)
    print("outer:", x)


outer()
\end{dhkcode}

예상 출력:

\begin{dhkzsh}
inner: 20
outer: 20
\end{dhkzsh}

이번에는 \zsh{inner}의 \zsh{x}가 외부 함수의 바인딩을 대상으로 한다.

따라서 \zsh{x = 20}은 \zsh{inner}에 새로운 지역 바인딩을 만드는 것이 아니라 \zsh{outer}의 \zsh{x} 바인딩을 변경한다.

개념적으로 다음과 같은 구조다.

\begin{dhkdiagram}[in=6]
node n1 (1,1) :: outer() 실행 환경
node n2 (2,2) :: Cell x
node n3 (2,3) :: int(10)
node n4 (3,3) :: inner의 nonlocal x

edge n1 -- n2 :: VH
edge n2 -> n3
edge n4 -> n2 :: HV
\end{dhkdiagram}

\zsh{inner()}에서 대입이 실행되면 다음과 같이 바뀐다.

\begin{dhkdiagram}[in=6]
node n1 (1,1) :: outer() 실행 환경
node n2 (2,2) :: Cell x
node n3 (2,3) :: int(20)
node n4 (3,3) :: inner의 nonlocal x

edge n1 -- n2 :: VH
edge n2 -> n3
edge n4 -> n2 :: HV
\end{dhkdiagram}

여기서 셀 객체 자체가 교체되는 것은 아니다.

기존 셀이 참조하는 객체가 변경되는 것이다.

이 차이를 다음 절에서 직접 확인해 보자.

\subsubsection{셀 객체의 동일성과 바인딩 변경}

이번 학습에서 가장 중요한 실험이다.

다음 코드를 살펴보자.

\begin{dhkcode}
def outer():
    x = 10

    def get():
        return x

    def set_value(value):
        nonlocal x
        x = value

    return get, set_value


get, set_value = outer()

get_cell = get.__closure__[0]
set_cell = set_value.__closure__[0]

print(get_cell is set_cell)

print(get_cell.cell_contents)

set_value(100)

print(get_cell.cell_contents)

print(get_cell is set_cell)
\end{dhkcode}

예상 출력:

\begin{dhkzsh}
True
10
100
True
\end{dhkzsh}

이 결과에는 서로 다른 두 종류의 동일성 관계가 포함되어 있다.

\begin{dhknote}
\begin{dhksmry}[centered]
첫 번째 동일성 - 셀 객체의 동일성
\end{dhksmry}
\end{dhknote}

\begin{dhkcode}
get_cell is set_cell
\end{dhkcode}

결과는 \zsh{True}다.

\zsh{get}과 \zsh{set_value}는 같은 \zsh{outer()} 호출에서 생성되었고, 동일한 외부 변수 \zsh{x}를 참조하기 때문이다.

두 함수는 하나의 셀 객체를 공유한다.

\begin{dhknote}
\begin{dhksmry}[centered]
두 번째 관계 - 셀이 참조하는 객체
\end{dhksmry}
\end{dhknote}

처음에는 다음과 같다.

\begin{dhkcode}
get_cell.cell_contents
\end{dhkcode}

결과:

\begin{dhkzsh}
10
\end{dhkzsh}

그런데 다음 호출을 실행하면

\begin{dhkcode}
set_value(100)
\end{dhkcode}

\zsh{set_value} 내부에서 다음 문장이 실행된다.

\begin{dhkcode}
x = value
\end{dhkcode}

\zsh{nonlocal x}가 선언되어 있으므로 외부 함수의 \zsh{x} 바인딩이 변경된다.

따라서 셀이 참조하는 객체가 정수 객체 10에서 정수 객체 100으로 바뀐다.

하지만 셀 객체는 그대로다.

\begin{dhknote}
\begin{dhksmry}[centered]
실행 전
\end{dhksmry}
\end{dhknote}

\begin{dhkdiagram}[in=6]
node n1 (1,1) :: get
node n2 (1,2) :: 함수 객체 1
node n3 (3,2) :: Cell x
node n4 (5,1) :: set_value
node n5 (5,2) :: 함수 객체 2
node n6 (3,3) :: int(10)

edge n1 -> n2
edge n2 -> n3
edge n4 -> n5
edge n5 -> n3
edge n3 -> n6
\end{dhkdiagram}

\begin{dhknote}
\begin{dhksmry}[centered]
실행 후
\end{dhksmry}
\end{dhknote}

\begin{dhkdiagram}[in=6]
node n1 (1,1) :: get
node n2 (1,2) :: 함수 객체 1
node n3 (3,2) :: Cell x
node n4 (5,1) :: set_value
node n5 (5,2) :: 함수 객체 2
node n6 (3,3) :: int(100)

edge n1 -> n2
edge n2 -> n3
edge n4 -> n5
edge n5 -> n3
edge n3 -> n6
\end{dhkdiagram}

두 그림에서 바뀐 것은 \zsh{Cell x}가 참조하는 객체다.

\zsh{get} 함수 객체도 그대로이고, \zsh{set_value} 함수 객체도 그대로이며, 셀 객체 역시 동일하다.

이를 다음처럼 정리할 수 있다.

\begin{dhktable}
|비교 대상|	set_value(100) 실행 전후|
|:---|:---|
|get 함수 객체|	동일|
|set_value 함수 객체|	동일|
|공유하는 셀 객체|	동일|
|셀이 참조하는 정수 객체|	변경|
|get() 반환값|	10 → 100|
\end{dhktable}

\zsh{nonlocal}에 의한 재바인딩은 셀 객체를 교체하는 것이 아니라, 기존 셀의 바인딩 대상을 변경하는 연산이다.

\subsubsection{\zsh{co_freevars}, \zsh{co_cellvars}, \zsh{__closure__}로 검증하기}

앞의 예제를 조금 확장해 보자.

\begin{dhkcode}
def outer():
    x = 10

    def get():
        return x

    def set_value(value):
        nonlocal x
        x = value

    return get, set_value


get, set_value = outer()

print("outer cellvars:", outer.__code__.co_cellvars)

print("get freevars:", get.__code__.co_freevars)

print("set freevars:", set_value.__code__.co_freevars)

print(
    "same cell:",
    get.__closure__[0] is set_value.__closure__[0]
)
\end{dhkcode}

예상 출력:

\begin{dhkzsh}
outer cellvars: ('x',)
get freevars: ('x',)
set freevars: ('x',)
same cell: True
\end{dhkzsh}

이 결과를 순서대로 해석해 보자.

\begin{dhknote}
\begin{dhksmry}[centered]
1. outer.__code__.co_cellvars
\end{dhksmry}
\end{dhknote}

\begin{dhkzsh}
('x',)
\end{dhkzsh}

\zsh{outer}의 지역 변수 \zsh{x}는 중첩 함수에서 접근해야 하는 바인딩이다.

따라서 셀 변수로 분류된다.

\begin{dhknote}
\begin{dhksmry}[centered]
2. get.__code__.co_freevars
\end{dhksmry}
\end{dhknote}

\begin{dhkzsh}
('x',)
\end{dhkzsh}

\zsh{get}은 외부 함수의 \zsh{x}를 읽어야 한다.

따라서 \zsh{x}는 \zsh{get}의 자유 변수다.

\begin{dhknote}
\begin{dhksmry}[centered]
3. __closure__ 비교
\end{dhksmry}
\end{dhknote}

두 함수의 \zsh{__closure__}에 포함된 셀 객체를 비교하면 동일한 객체라는 사실을 확인할 수 있다.

여기서 중요한 결론은 다음과 같다.

\zsh{nonlocal}이 선언되었다고 해서 해당 이름이 현재 함수의 지역 변수가 되는 것이 아니다.

오히려 \zsh{nonlocal}은 그 이름을 현재 함수의 지역 바인딩 대상으로 취급하지 않도록 지정한다.

따라서 \zsh{set_value}의 \zsh{x}는 외부 함수의 바인딩을 사용하는 자유 변수로 분류된다.

\subsubsection{객체 변경(Mutation)과 이름 재바인딩(Rebinding)}

이제 \zsh{nonlocal}을 사용할 때 가장 자주 혼동하는 문제를 살펴보자.

다음 두 표현식은 서로 다른 연산이다.

\begin{dhkcode}
values.append(10)

values = [10]
\end{dhkcode}

첫 번째는 \zsh{values}가 참조하는 리스트 객체의 내부 상태를 변경한다.

두 번째는 \zsh{values}라는 이름을 새로운 리스트 객체에 바인딩한다.

이 차이를 클로저에 적용해 보자.

\begin{dhknote}
\begin{dhksmry}[centered]
사례 1. 리스트 객체의 상태 변경
\end{dhksmry}
\end{dhknote}

\begin{dhkcode}
def outer():
    values = []

    def add(value):
        values.append(value)
        return values

    return add


add = outer()

print(add(10))
print(add(20))
print(add(30))
\end{dhkcode}

예상 출력:

\begin{dhkzsh}
[10]
[10, 20]
[10, 20, 30]
\end{dhkzsh}

여기에는 \zsh{nonlocal} 선언이 없다.

그런데도 외부 함수에서 생성한 리스트의 상태가 변경되고 있다.

왜 가능할까?

\zsh{add()} 내부에서는 \zsh{values}라는 이름에 새로운 객체를 대입하지 않기 때문이다.

다음 표현식을 분석해 보자.

\begin{dhkcode}
values.append(value)
\end{dhkcode}

실행 과정은 다음과 같다.

\begin{dhksmry}
1. 자유 변수 values를 통해 리스트 객체를 조회한다.

2. 그 리스트 객체에서 append 메서드를 조회한다.

3. 현재 호출의 매개변수 value가 참조하는 객체를 인수로 전달한다.

4. append 메서드가 기존 리스트 객체를 변경한다.
\end{dhksmry}

이 과정에서 외부 이름 \zsh{values}의 바인딩 자체는 변경되지 않는다.

따라서 \zsh{nonlocal values}가 필요하지 않은 것이다.

\begin{dhknote}
\begin{dhksmry}[centered]
객체 구조
\end{dhksmry}
\end{dhknote}

\begin{dhkdiagram}[in=6]
node n1 (1,1) :: add 함수 객체
node n2 (3,1) :: Cell values
node n3 (5,1) :: List 객체
node n4 (5,2) :: [10, 20, 30]

edge n1 -> n2
edge n2 -> n3
edge n3 -> n4
\end{dhkdiagram}

리스트 객체는 계속 동일하다.

그 내부 원소 구성이 변경되는 것이다.

\subsubsection{리스트를 재바인딩하려면?}

이번에는 다음 코드를 살펴보자.

\begin{dhkcode}
def outer():
    values = []

    def reset():
        values = []
        return values

    def get():
        return values

    return reset, get


reset, get = outer()

print(get())
print(reset())
print(get())
\end{dhkcode}

예상 출력:

\begin{dhkzsh}
[]
[]
[]
\end{dhkzsh}

세 출력이 모두 빈 리스트이므로 얼핏 보면 \zsh{reset()}이 외부 리스트를 초기화한 것처럼 보일 수 있다.

하지만 실제로는 그렇지 않다.

\zsh{reset()} 내부의 다음 문장에 주목해 보자.

\begin{dhkcode}
values = []
\end{dhkcode}

\zsh{nonlocal} 선언이 없으므로 \zsh{values}는 \zsh{reset}의 지역 이름이다.

따라서 \zsh{reset()}은 새로운 리스트 객체를 생성하고 자신의 지역 이름에 바인딩할 뿐이다.

외부 함수의 \zsh{values}에는 영향을 주지 않는다.

이를 명확하게 확인해 보자.

\begin{dhkcode}
def outer():
    values = [10, 20]

    def reset():
        values = []
        return values

    def get():
        return values

    return reset, get


reset, get = outer()

print("before:", get())
print("reset result:", reset())
print("after:", get())
\end{dhkcode}

예상 출력:

\begin{dhkzsh}
before: [10, 20]
reset result: []
after: [10, 20]
\end{dhkzsh}

이제 차이가 분명해졌다.

\zsh{reset()}에서 새로운 리스트 객체를 생성했지만, 외부 함수의 \zsh{values}는 변경되지 않았다.

\begin{dhknote}
\begin{dhksmry}[centered]
nonlocal을 사용한 경우
\end{dhksmry}
\end{dhknote}

\begin{dhkcode}
def outer():
    values = [10, 20]

    def reset():
        nonlocal values
        values = []
        return values

    def get():
        return values

    return reset, get


reset, get = outer()

print("before:", get())
print("reset result:", reset())
print("after:", get())
\end{dhkcode}

예상 출력:

\begin{dhkzsh}
before: [10, 20]
reset result: []
after: []
\end{dhkzsh}

이번에는 \zsh{reset()}이 외부 함수의 \zsh{values} 바인딩을 새로운 리스트 객체로 변경한다.

\begin{dhknote}
\begin{dhksmry}[centered]
실행 전
\end{dhksmry}
\end{dhknote}

\begin{dhksmry}[centered]
공유 Cell values
\da
List 객체 A
[10, 20]
\end{dhksmry}

\begin{dhknote}
\begin{dhksmry}[centered]
실행 후
\end{dhksmry}
\end{dhknote}

\begin{dhksmry}[centered]
공유 Cell values
\da
List 객체 B
[]
\end{dhksmry}

여기서 기존 리스트 객체 A가 반드시 즉시 소멸한다는 뜻은 아니다.

다른 이름이나 객체가 리스트 A를 참조하고 있다면 계속 존재할 수 있다.

핵심은 셀의 참조 대상이 리스트 A에서 리스트 B로 변경되었다는 사실이다.

\subsubsection{같은 결과처럼 보여도 객체 관계는 다를 수 있다}

이번에는 \zsh{clear()}와 재바인딩을 비교해 보자.

\begin{dhkcode}
def make_state():
    values = [10, 20]

    def get():
        return values

    def clear_in_place():
        values.clear()

    def replace():
        nonlocal values
        values = []

    return get, clear_in_place, replace


get, clear_in_place, replace = make_state()

original = get()

print("initial:", original)

clear_in_place()

print("after clear:", original)

replace()

print("after replace:", original)
print("current:", get())

print(original is get())
\end{dhkcode}

예상 출력:

\begin{dhkzsh}
initial: [10, 20]
after clear: []
after replace: []
current: []
False
\end{dhkzsh}

이 결과는 특히 중요하다.

\zsh{clear_in_place()}와 \zsh{replace()}를 실행한 뒤 모두 빈 리스트가 나타나지만, 두 함수가 수행한 연산은 다르다.

\begin{dhknote}
\begin{dhksmry}[centered]
clear_in_place()
\end{dhksmry}
\end{dhknote}

\begin{dhkcode}
values.clear()
\end{dhkcode}

기존 리스트 객체의 내부 상태를 변경한다.

따라서 \zsh{original}이 참조하는 리스트 객체도 비워진다.

\begin{dhknote}
\begin{dhksmry}[centered]
replace()
\end{dhksmry}
\end{dhknote}

\begin{dhkcode}
nonlocal values
values = []
\end{dhkcode}

새로운 리스트 객체를 생성하고 외부 바인딩을 변경한다.

따라서 \zsh{original}은 여전히 기존 리스트 객체를 참조하지만, \zsh{get()}은 새로운 리스트 객체를 반환한다.

마지막 비교가 이를 보여준다.

\begin{dhkcode}
original is get()
\end{dhkcode}

결과는 \zsh{False}다.

이 차이는 클래스 인스턴스에서도 동일하게 나타난다.

\begin{dhkcode}
self.values.clear()
\end{dhkcode}

와

\begin{dhkcode}
self.values = []
\end{dhkcode}

는 서로 다른 연산이다.

전자는 기존 리스트 객체를 변경하고, 후자는 일반적인 속성 대입 규칙에 따라 \zsh{values} 속성의 참조 대상을 변경한다.

따라서 클로저의 상태 변경 원리는 기존에 학습한 객체 모델과 분리된 특별한 규칙이 아니다.

Python의 일반적인 객체 변경과 바인딩 변경 원리가 클로저 환경에서도 적용되는 것이다.

\subsubsection{\zsh{+=} 연산은 왜 주의해야 할까?}

이번에는 다음 코드를 살펴보자.

\begin{dhkcode}
def outer():
    count = 0

    def increment():
        count += 1
        return count

    return increment


increment = outer()

print(increment())
\end{dhkcode}

예상 결과:

\begin{dhkzsh}
UnboundLocalError
\end{dhkzsh}

왜 이런 결과가 발생할까?

\zsh{count += 1}은 단순히 외부 정수 객체의 상태를 변경하는 연산이 아니다.

정수 객체는 불변 객체(Immutable Object)다.

또한 \zsh{count += 1}은 이름 \zsh{count}에 대한 대입을 포함하는 증강 대입문(Augmented Assignment)이다.

따라서 별도의 \zsh{nonlocal} 선언이 없으면 \zsh{count}는 \zsh{increment}의 지역 이름으로 분류된다.

그런데 증강 대입에서는 기존 값을 먼저 읽어야 한다.

현재 지역 \zsh{count}에는 아직 객체가 바인딩되지 않았으므로 

\zsh{UnboundLocalError}가 발생하는 것이다.

\begin{dhknote}
\begin{dhksmry}[centered]
해결 방법
\end{dhksmry}
\end{dhknote}

\begin{dhkcode}
def outer():
    count = 0

    def increment():
        nonlocal count
        count += 1
        return count

    return increment


increment = outer()

print(increment())
print(increment())
print(increment())
\end{dhkcode}

예상 출력:

\begin{dhkzsh}
1
2
3
\end{dhkzsh}

이번에는 \zsh{count}가 외부 함수의 바인딩을 대상으로 한다.

따라서 기존 정수 객체를 읽고, 덧셈 결과를 다시 외부 이름에 바인딩할 수 있다.

\begin{dhknote}
\begin{dhksmry}[centered]
실행 과정
\end{dhksmry}
\end{dhknote}

첫 번째 호출:

\begin{dhksmry}
    Cell count \ra int(0)
\end{dhksmry}

\zsh{count += 1} 실행 후:

\begin{dhksmry}
    Cell count \ra int(1)
\end{dhksmry}

두 번째 호출 후:

\begin{dhksmry}
    Cell count \ra int(2)
\end{dhksmry}

세 번째 호출 후:

\begin{dhksmry}
    Cell count \ra int(3)
\end{dhksmry}

여기서 정수 객체 자체가 변경되는 것은 아니다.

매번 새로운 계산 결과에 바인딩되는 것이다.

\subsubsection{가변 객체의 \zsh{+=}는 어떻게 다를까?}

이제 정수 대신 리스트를 사용해 보자.

\begin{dhkcode}
def outer():
    values = []

    def add(value):
        values += [value]
        return values

    return add


add = outer()

print(add(10))
\end{dhkcode}

이 코드는 어떻게 될까?

리스트는 가변 객체이므로 정상적으로 동작할 것 같지만,

실제로는 \zsh{UnboundLocalError}가 발생한다.

왜 그럴까?

여기서 중요한 것은 객체가 가변인지 여부와 이름의 스코프 판정은 별개의 문제라는 사실이다.

\zsh{values += [value]}는 증강 대입문이므로 \zsh{values}를 바인딩 대상으로 사용한다.

따라서 \zsh{nonlocal} 선언이 없다면 \zsh{values}는 현재 함수의 지역 이름으로 판정된다.

그런데 기존 지역 \zsh{values}가 아직 바인딩되지 않았으므로 읽는 과정에서 오류가 발생하는 것이다.

\begin{dhknote}
\begin{dhksmry}[centered]
nonlocal을 추가하면?
\end{dhksmry}
\end{dhknote}

\begin{dhkcode}
def outer():
    values = []

    def add(value):
        nonlocal values
        values += [value]
        return values

    return add


add = outer()

print(add(10))
print(add(20))
\end{dhkcode}

예상 출력:

\begin{dhkzsh}
[10]
[10, 20]
\end{dhkzsh}

이제 정상적으로 동작한다.

그런데 여기서 한 단계 더 생각해야 한다.

리스트의 \zsh{+=}는 일반적으로 기존 리스트 객체를 직접 변경하는 

\zsh{__iadd__} 동작을 사용할 수 있다.

따라서 \zsh{nonlocal values}가 있다고 해서 반드시 새로운 리스트 객체가 생성되는 것은 아니다.

다음 코드로 확인해 보자.

\begin{dhkcode}
def outer():
    values = []

    def add(value):
        nonlocal values

        before = values

        values += [value]

        return before is values

    return add


add = outer()

print(add(10))
print(add(20))
\end{dhkcode}

일반적인 Python 리스트에서 예상 출력:

\begin{dhkzsh}
True
True
\end{dhkzsh}

이 결과는 다음 두 가지를 동시에 보여준다.

\begin{dhksmry}
    1. nonlocal은 이름의 바인딩 대상을 지정한다.

    2. 시제 증강 대입이 기존 객체를 변경할지 새로운 객체를 반환할지는
       피연산자의 타입과 연산 구현에 따라 달라질 수 있다.
\end{dhksmry}

따라서 다음 세 연산을 동일하게 취급하면 안 된다.

\begin{dhkcode}
values.append(10)

values += [10]

values = values + [10]
\end{dhkcode}

일반적인 리스트에서는 첫 번째와 두 번째가 기존 리스트 객체를 변경하지만, 세 번째는 새로운 리스트 객체를 만들어 이름을 재바인딩한다.

또한 두 번째와 세 번째는 이름 대입을 포함하므로 외부 함수의 \zsh{values}를 대상으로 하려면 \zsh{nonlocal}이 필요하다.

객체 변경 여부와 이름의 스코프 판정 여부는 서로 다른 기준으로 판단해야 한다.

\subsubsection{\zsh{nonlocal}은 가장 가까운 외부 함수의 바인딩을 찾는다}

지금까지는 외부 함수가 한 단계만 존재했다.

이번에는 세 단계의 중첩 함수를 살펴보자.

\begin{dhkcode}
def level1():
    x = 10

    def level2():
        x = 20

        def level3():
            nonlocal x
            x = 30

        level3()

        print("level2:", x)

    level2()

    print("level1:", x)


level1()
\end{dhkcode}

예상 출력:

\begin{dhkzsh}
level2: 30
level1: 10
\end{dhkzsh}

왜 \zsh{level1}의 \zsh{x}는 변경되지 않았을까?

\zsh{level3}의 렉시컬 외부 함수 스코프를 살펴보자.

\begin{dhkdir}[in=6]
level1
> x = 10
> level2
>> x = 20
>> level3
>>> nonlocal x
\end{dhkdir}

\zsh{level3}에서 \zsh{nonlocal x}를 선언하면 현재 함수의 지역 스코프를 제외하고, 렉시컬 외부 함수 스코프에서 \zsh{x}의 기존 바인딩을 찾는다.

가장 가까운 외부 함수인 \zsh{level2}에 \zsh{x}가 존재하므로 그 바인딩을 사용한다.

따라서 \zsh{level2}의 \zsh{x}가 30으로 변경된다.

반면 \zsh{level1}의 \zsh{x}는 그대로 10이다.

\begin{dhknote}
\begin{dhksmry}[centered]
주의할 점
\end{dhksmry}
\end{dhknote}

\zsh{nonlocal}은 특정 외부 함수를 이름으로 지정하는 문법을 제공하지 않는다.

따라서 다음처럼 작성할 수는 없다.

\begin{dhkcode}
nonlocal level1.x
\end{dhkcode}

이는 유효한 Python 문법이 아니다.

\zsh{nonlocal}은 점 표기법을 통한 속성 탐색이 아니라 렉시컬 외부 함수 스코프의 이름 바인딩을 선택하는 선언이기 때문이다.

\subsubsection{중간 함수에서 이름을 바인딩하지 않는 경우}

다음 코드를 살펴보자.

\begin{dhkcode}
def level1():
    x = 10

    def level2():

        def level3():
            nonlocal x
            x = 30

        level3()

    level2()

    return x


print(level1())
\end{dhkcode}

예상 출력:

\begin{dhkzsh}
30
\end{dhkzsh}

이번에는 \zsh{level2}에 \zsh{x} 바인딩이 없다.

따라서 \zsh{level3}의 \zsh{nonlocal x}는 더 바깥쪽 함수인 \zsh{level1}의 바인딩을 대상으로 한다.

여기서 중요한 사실이 있다.

\zsh{nonlocal}은 반드시 바로 바깥 함수에만 접근하는 선언이 아니다.

렉시컬 외부 함수 스코프들을 탐색하여 가장 가까운 적절한 기존 바인딩을 사용한다.

이 원리는 \zsh{A.3}에서 학습한 Enclosing Scope와 연결된다.

\subsubsection{\zsh{nonlocal}의 제약과 예외}

\zsh{nonlocal}은 아무 이름에나 사용할 수 있는 선언이 아니다.

대표적인 제약을 살펴보자.

\begin{dhknote}
\begin{dhksmry}[centered]
1. 외부 함수 스코프에 바인딩이 존재해야 한다
\end{dhksmry}
\end{dhknote}

\begin{dhkcode}
def outer():

    def inner():
        nonlocal x
        x = 10

    return inner
\end{dhkcode}

이 코드는 \zsh{SyntaxError}를 발생시킨다.

왜냐하면 \zsh{inner}를 감싸는 외부 함수 스코프에 \zsh{x}라는 바인딩이 존재하지 않기 때문이다.

\zsh{nonlocal}은 새로운 외부 바인딩을 생성하기 위한 선언이 아니다.

기존 렉시컬 외부 함수 바인딩을 대상으로 하는 선언이다.

\begin{dhknote}
\begin{dhksmry}[centered]
2. 모듈 전역 이름만으로는 충분하지 않다
\end{dhksmry}
\end{dhknote}

\begin{dhkcode}
x = 10

def func():
    nonlocal x
    x = 20
\end{dhkcode}

이 코드 역시 \zsh{SyntaxError}를 발생시킨다.

모듈 전역의 \zsh{x}는 \zsh{nonlocal}의 대상이 아니기 때문이다.

전역 이름을 재바인딩하려면 \zsh{global}을 사용해야 한다.

\begin{dhkcode}
x = 10

def func():
    global x
    x = 20


func()

print(x)
\end{dhkcode}

예상 출력:

\begin{dhkzsh}
20
\end{dhkzsh}

\begin{dhknote}
\begin{dhksmry}[centered]
3. 함수 내부에서 사용한 뒤 선언할 수 없다
\end{dhksmry}
\end{dhknote}

\begin{dhkcode}
def outer():
    x = 10

    def inner():
        print(x)
        nonlocal x
\end{dhkcode}

이 코드는 \zsh{SyntaxError}를 발생시킨다.

같은 코드 블록에서 \zsh{nonlocal x} 선언보다 앞서 \zsh{x}를 사용했기 때문이다.

\zsh{nonlocal}은 일반적인 실행 명령이 아니라 해당 코드 블록의 이름 분류에 영향을 주는 선언이므로, 이러한 제약이 존재한다.

\begin{dhknote}
\begin{dhksmry}[centered]
4. 외부 바인딩이 아직 객체를 참조하지 않을 수도 있다
\end{dhksmry}
\end{dhknote}

다음 코드를 보자.

\begin{dhkcode}
def outer():

    def inner():
        nonlocal x
        return x

    result = inner()

    x = 10

    return result


outer()
\end{dhkcode}

이 코드는 컴파일 자체는 가능하다.

왜냐하면 \zsh{outer} 내부에 \zsh{x = 10}이라는 바인딩 구문이 존재하기 때문이다.

하지만 실제 실행에서는 \zsh{NameError}가 발생한다.

\zsh{inner()}를 호출하는 시점에 \zsh{outer}의 \zsh{x = 10}이 아직 실행되지 않았기 때문이다.

이 차이는 \zsh{A.3}에서 배운 내용과 동일하다.

이름의 정적 분류가 가능하다는 사실과 실행 시점에 객체가 바인딩되어 있다는 사실은 서로 다르다.

\subsubsection{\zsh{global}과 \zsh{nonlocal}의 객체 모델 비교}

이번에는 두 선언을 같은 코드에서 비교해 보자.

\begin{dhkcode}
x = 10


def outer():
    x = 20

    def change_global():
        global x
        x = 100

    def change_nonlocal():
        nonlocal x
        x = 200

    change_global()
    change_nonlocal()

    return x


result = outer()

print("global:", x)
print("outer:", result)
\end{dhkcode}

예상 출력:

\begin{dhkzsh}
global: 100
outer: 200
\end{dhkzsh}

두 선언은 서로 다른 바인딩을 대상으로 한다.

\begin{dhknote}
\begin{dhksmry}[centered]
global x
\end{dhksmry}
\end{dhknote}

\begin{dhkdir}[in=6]
모듈 전역 네임스페이스
> x -> int(100)
\end{dhkdir}

\begin{dhknote}
\begin{dhksmry}[centered]
nonlocal x
\end{dhksmry}
\end{dhknote}

\begin{dhkdir}[in=6]
outer() 실행 환경
> Cell x -> int(200)
\end{dhkdir}

여기서 다음 구분이 중요하다.

\begin{dhktable}
|비교 기준|	global|	nonlocal|
|:---|:---|:---|
|대상|	모듈 전역 바인딩|	렉시컬 외부 함수 바인딩|
|기존 바인딩 필요|	선언 시 불필요|	외부 함수 스코프에 바인딩 구문 필요|
|대표적인 접근 구조|	전역 네임스페이스|	클로저 셀|
|선언의 목적|	전역 이름의 바인딩 지정|	외부 함수 이름의 바인딩 지정|
|객체 자체를 변경하는가?|	선언 자체는 변경하지 않음|	선언 자체는 변경하지 않음|
\end{dhktable}

또한 \zsh{global}과 \zsh{nonlocal}은 서로 대체 가능한 문법이 아니다.

어떤 스코프의 이름을 변경하려는지에 따라 선택해야 한다.

\subsubsection{클로저를 이용한 상태 관리 - Counter}

이제 지금까지 학습한 개념을 하나의 함수 생성기에 적용해 보자.

\begin{dhkcode}
def make_counter(initial=0):
    count = initial

    def increment(step=1):
        nonlocal count

        count += step

        return count

    def get():
        return count

    def reset():
        nonlocal count

        count = initial

    return increment, get, reset


increment, get, reset = make_counter(10)

print(get())

print(increment())
print(increment(5))

reset()

print(get())
\end{dhkcode}

예상 출력:

\begin{dhkzsh}
10
11
16
10
\end{dhkzsh}

이 코드를 객체 관계로 분석해 보자.

\begin{dhknote}
\begin{dhksmry}[centered]
1. make_counter(10) 호출
\end{dhksmry}
\end{dhknote}

매개변수 \zsh{initial}에 정수 객체 10이 바인딩된다.

다음 문장이 실행된다.

\begin{dhkcode}
count = initial
\end{dhkcode}

\zsh{count} 역시 정수 객체 10을 참조한다.

여기서 \zsh{initial}과 \zsh{count}는 서로 다른 이름이지만 처음에는 동일한 정수 객체를 참조한다.

\begin{dhknote}
\begin{dhksmry}[centered]
2. 내부 함수 객체 생성
\end{dhksmry}
\end{dhknote}

세 개의 함수 객체가 생성된다.

\begin{dhkcode}
increment
get
reset
\end{dhkcode}

각 함수는 필요한 외부 바인딩에 접근하기 위한 클로저 환경을 가진다.

여기서 중요한 점이 있다.

\zsh{increment}와 \zsh{get}은 \zsh{count}를 참조한다.

\zsh{reset}은 \zsh{count}와 \zsh{initial}을 참조한다.

따라서 \zsh{count}에 대응하는 셀은 세 함수가 공유한다.

반면 \zsh{initial}에 대응하는 셀은 \zsh{reset}에서만 필요하다.

이를 구조적으로 나타내면 다음과 같다.

\begin{dhkdiagram}[in=6]
node n1 (1,1) :: increment 함수 객체
node n2 (2,1) :: get 함수 객체
node n3 (3,1) :: reset 함수 객체
node n4 (2,2) :: Cell count
node n5 (2,3) :: int(10)
node n6 (4,2) :: Cell initial
node n7 (4,3) :: int(10)

edge n1 -> n4 :: HV
edge n2 -> n4
edge n3 -> n4 :: HV
edge n4 -> n5
edge n3 -> n6 :: VH
edge n6 -> n7
\end{dhkdiagram}

\begin{dhknote}
\begin{dhksmry}[centered]
3. increment() 호출
\end{dhksmry}
\end{dhknote}

\begin{dhkcode}
count += step
\end{dhkcode}

\zsh{count}는 외부 함수의 바인딩이다.

\zsh{nonlocal count} 선언에 따라 기존 값을 읽고 새로운 계산 결과를 다시 그 바인딩에 연결한다.

따라서 \zsh{count}는 11을 참조하게 된다.

\begin{dhknote}
\begin{dhksmry}[centered]
4. increment(5) 호출
\end{dhksmry}
\end{dhknote}

이번에는 \zsh{step}이 5를 참조한다.

기존 \zsh{count}는 11이므로 계산 결과는 16이다.

이후 \zsh{count}가 16에 재바인딩된다.

\begin{dhknote}
\begin{dhksmry}[centered]
5. reset() 호출
\end{dhksmry}
\end{dhknote}

\begin{dhkcode}
count = initial
\end{dhkcode}

\zsh{count}는 \zsh{nonlocal} 선언에 의해 외부 바인딩을 대상으로 한다.

\zsh{initial}은 외부 함수의 자유 변수로 읽힌다.

따라서 \zsh{count}가 다시 정수 객체 10을 참조하게 된다.

여기서 새로운 \zsh{make_counter()} 호출이 발생하는 것은 아니다.

기존 클로저 환경의 바인딩을 변경하는 것이다.

\subsubsection{\zsh{__closure__}를 이용한 상태 공유 검증}

앞의 예제에서 실제로 같은 셀을 공유하는지 확인해 보자.

\begin{dhkcode}
def make_counter(initial=0):
    count = initial

    def increment(step=1):
        nonlocal count
        count += step
        return count

    def get():
        return count

    def reset():
        nonlocal count
        count = initial

    return increment, get, reset


increment, get, reset = make_counter(10)


def closure_map(func):
    if func.__closure__ is None:
        return {}

    return dict(
        zip(
            func.__code__.co_freevars,
            func.__closure__
        )
    )


increment_cells = closure_map(increment)
get_cells = closure_map(get)
reset_cells = closure_map(reset)

print(increment_cells.keys())
print(get_cells.keys())
print(reset_cells.keys())

print(
    increment_cells["count"]
    is get_cells["count"]
)

print(
    get_cells["count"]
    is reset_cells["count"]
)

print(
    reset_cells["initial"].cell_contents
)
\end{dhkcode}

예상 출력:

\begin{dhkzsh}
dict_keys(['count'])
dict_keys(['count'])
dict_keys(['count', 'initial'])
True
True
10
\end{dhkzsh}

이 예제에서 \zsh{closure_map()}은 함수 객체의 자유 변수 이름과 셀 객체를 대응시키는 보조 함수다.

여기서 다음 사실을 직접 확인할 수 있다.

\begin{dhksmry}
    - increment, get, reset은 모두 count 셀을 공유한다.

    - reset은 추가로 initial 셀을 참조한다.

    - initial과 count는 처음에 같은 정수 객체를 참조하더라도
      서로 다른 바인딩이다.
\end{dhksmry}

이제 다음 비교를 생각해 보자.

\begin{dhkcode}
reset_cells["initial"] is reset_cells["count"]
\end{dhkcode}

예상 결과:

\begin{dhkzsh}
False
\end{dhkzsh}

왜냐하면 두 이름은 서로 다른 바인딩이기 때문이다.

처음에 동일한 정수 객체를 참조한다는 사실만으로 같은 셀 객체를 사용해야 하는 것은 아니다.

이것은 \zsh{A.1}에서 배운 객체의 동일성과 이름 바인딩의 동일성을 구분해야한다는 원리가 클로저에도 적용된 사례다.

\subsubsection{서로 다른 Counter 인스턴스의 상태 독립성}

이번에는 함수 생성기를 두 번 호출해 보자.

\begin{dhkcode}
def make_counter():
    count = 0

    def increment():
        nonlocal count
        count += 1
        return count

    return increment


counter_a = make_counter()
counter_b = make_counter()

print(counter_a())
print(counter_a())

print(counter_b())

print(
    counter_a.__closure__[0]
    is counter_b.__closure__[0]
)
\end{dhkcode}

예상 출력:

\begin{dhkzsh}
1
2
1
False
\end{dhkzsh}

여기서 \zsh{counter_a}와 \zsh{counter_b}는 서로 다른 함수 객체다.

그리고 각각 다른 \zsh{make_counter()} 호출에서 생성되었다.

따라서 서로 독립적인 \zsh{count} 바인딩을 가진다.

\begin{dhkdiagram}[in=6]
node n1 (1,1) :: counter_a
node n2 (3,1) :: 함수 객체 A
node n3 (5,1) :: Cell count A
node n4 (5,2) :: int(2)

edge n1 -> n2
edge n2 -> n3
edge n3 -> n4
\end{dhkdiagram}

\begin{dhkdiagram}[in=6]
node n1 (1,1) :: counter_b
node n2 (3,1) :: 함수 객체 B
node n3 (5,1) :: Cell count B
node n4 (5,2) :: int(1)

edge n1 -> n2
edge n2 -> n3
edge n3 -> n4
\end{dhkdiagram}

이 구조는 클래스의 서로 다른 인스턴스가 독립적인 상태를 가지는 것과 유사하다.

다만 상태가 저장되는 위치는 다르다.

클래스에서는 일반적으로 인스턴스 속성을 통해 상태를 관리하지만, 여기서는 함수 객체가 참조하는 클로저 셀을 통해 상태를 관리하는 것이다.

\subsubsection{클래스 인스턴스와 클로저의 상태 관리 비교}

다음 두 구현을 비교해 보자.

\begin{dhknote}
\begin{dhksmry}[centered]
클로저 기반 구현
\end{dhksmry}
\end{dhknote}

\begin{dhkcode}
def make_counter():
    count = 0

    def increment():
        nonlocal count
        count += 1
        return count

    return increment
\end{dhkcode}

\begin{dhknote}
\begin{dhksmry}[centered]
클래스 기반 구현
\end{dhksmry}
\end{dhknote}

\begin{dhkcode}
class Counter:
    def __init__(self):
        self.count = 0

    def increment(self):
        self.count += 1
        return self.count
\end{dhkcode}

두 구현 모두 호출 사이에 상태를 유지할 수 있다.

하지만 상태에 접근하는 메커니즘은 다르다.

\begin{dhknote}
\begin{dhksmry}[centered]
클로저 기반
\end{dhksmry}
\end{dhknote}

\begin{dhksmry}[centered]
increment 함수 객체
\da
__closure__
\da
Cell count
\da
int(0)
\end{dhksmry}

\begin{dhknote}
\begin{dhksmry}[centered]
클래스 기반
\end{dhksmry}
\end{dhknote}

\begin{dhkdiagram}[in=6]
node n1 (1,1) :: Counter 인스턴스
node n2 (3,1) :: __dict__
node n3 (5,2) :: "count"
node n4 (5,3) :: int(0)

edge n1 -> n2
edge n2 -- n3 ::VH
edge n3 -> n4
\end{dhkdiagram}

일반적인 사용자 정의 클래스에서 \zsh{self.count}는 인스턴스 속성 탐색과 대입을 통해 관리된다.

반면 클로저의 \zsh{count}는 자유 변수 바인딩을 통해 접근한다.

\begin{dhktable}
|비교 기준|	클로저|	클래스|
|:---|:---|:---|
|상태 접근|	자유 변수|	인스턴스 속성|
|상태 저장|	셀과 참조 객체|	인스턴스 네임스페이스 등|
|상태 재바인딩|	nonlocal|	self.attribute = ...|
|여러 동작의 상태 공유|	같은 셀을 공유하는 함수들|	같은 인스턴스를 참조하는 메서드들|
|독립적인 상태 생성|	외부 함수 재호출|	인스턴스 생성|
|대표적인 활용|	함수 생성기, 데코레이터|	명시적인 상태 객체, 복잡한 인터페이스|
\end{dhktable}

여기서 중요한 것은 클로저가 클래스를 대체하는 문법이라는 뜻이 아니다.

둘은 서로 다른 객체 구조를 통해 상태를 관리하는 방법이다.

문제의 규모와 필요한 인터페이스에 따라 적합한 방식이 달라진다.

\subsubsection{실전 응용 - 데코레이터의 호출 횟수 관리}

이제 클로저 상태 변경을 데코레이터와 연결해 보자.

다음 코드를 살펴보자.

\begin{dhkcode}
def count_calls(func):
    count = 0

    def wrapper(*args, **kwargs):
        nonlocal count

        count += 1

        print(f"호출 횟수: {count}")

        return func(*args, **kwargs)

    return wrapper


def add(a, b):
    return a + b


add = count_calls(add)

print(add(10, 20))
print(add(30, 40))
\end{dhkcode}

예상 출력:

\begin{dhkzsh}
호출 횟수: 1
30
호출 횟수: 2
70
\end{dhkzsh}

이 코드는 이후 \zsh{A.6}과 \zsh{A.7}에서 자세히 학습할 함수 기반 데코레이터의 기본 구조다.

이번에는 클로저 상태에 집중해서 분석해 보자.

\begin{dhknote}
\begin{dhksmry}[centered]
1. 원본 함수 객체 생성
\end{dhksmry}
\end{dhknote}

\begin{dhkcode}
def add(a, b):
    return a + b
\end{dhkcode}

원본 \zsh{add} 함수 객체가 생성된다.

\begin{dhknote}
\begin{dhksmry}[centered]
2. count_calls(add) 호출
\end{dhksmry}
\end{dhknote}

\begin{dhkcode}
add = count_calls(add)
\end{dhkcode}

원본 함수 객체가 \zsh{count_calls}의 매개변수 \zsh{func}에 바인딩된다.

그리고 다음 문장이 실행된다.

\begin{dhkcode}
count = 0
\end{dhkcode}

호출 횟수를 나타내는 지역 바인딩이 생성된다.

\begin{dhknote}
\begin{dhksmry}[centered]
3. wrapper 함수 객체 생성
\end{dhksmry}
\end{dhknote}

\begin{dhkcode}
def wrapper(*args, **kwargs):
    nonlocal count

    count += 1

    print(f"호출 횟수: {count}")

    return func(*args, **kwargs)
\end{dhkcode}

\zsh{wrapper}는 외부 함수의 두 이름을 사용한다.

\begin{dhksmry}
    - count

    - func
\end{dhksmry}

따라서 이 두 이름은 \zsh{wrapper}의 자유 변수다.

\zsh{count}는 \zsh{nonlocal} 선언을 통해 재바인딩 대상이 되고,

\zsh{func}는 원본 함수 객체를 호출하기 위해 읽는 이름이다.

\begin{dhknote}
\begin{dhksmry}[centered]
4. wrapper 함수 객체 반환
\end{dhksmry}
\end{dhknote}

\zsh{count_calls()}는 \zsh{wrapper} 함수 객체를 반환한다.

그리고 모듈 이름 \zsh{add}가 반환된 함수 객체에 재바인딩된다.

이후 \zsh{count_calls()} 호출은 종료된다.

하지만 반환된 \zsh{wrapper} 함수 객체는 클로저를 통해 \zsh{count}와 \zsh{func}의 바인딩에 접근할 수 있다.

\begin{dhknote}
\begin{dhksmry}[centered]
객체 구조
\end{dhksmry}
\end{dhknote}

\begin{dhkbranch}[in=8]
node main :: 모듈 이름 add
arrow main
node main :: wrapper 함수 객체
arrow main
node main :: __closure__
branchline main -> b1 :: Cell count -> int(0)
handoff main -> b2 :: Cell func -> 원본 add 함수 객체
stop b2
\end{dhkbranch}

\begin{dhknote}
\begin{dhksmry}[centered]
5. 첫 번째 add(10, 20) 호출
\end{dhksmry}
\end{dhknote}

현재 모듈 이름 \zsh{add}는 \zsh{wrapper} 함수 객체를 참조한다.

따라서 \zsh{wrapper}가 실행된다.

\begin{dhkcode}
count += 1
\end{dhkcode}

\zsh{count}가 1로 재바인딩된다.

이후 원본 함수 객체를 호출한다.

\begin{dhkcode}
func(*args, **kwargs)
\end{dhkcode}

원본 \zsh{add}는 30을 반환한다.

그 결과 \zsh{wrapper}도 30을 반환한다.

\begin{dhknote}
\begin{dhksmry}[centered]
6. 두 번째 add(30, 40) 호출
\end{dhksmry}
\end{dhknote}

이번에도 같은 \zsh{wrapper} 함수 객체가 호출된다.

따라서 같은 클로저 셀을 사용한다.

기존 \zsh{count}가 1이므로 \zsh{count += 1} 실행 후 2가 된다.

이후 원본 함수 객체를 호출하고 70을 반환한다.

여기서 핵심은 다음과 같다.

\begin{dhksmry}[centered]
데코레이터가 호출 횟수를 유지하는 이유는 함수 객체가

호출 사이에 동일한 클로저 셀을 계속 참조하기 때문이다.
\end{dhksmry}

\subsubsection{\zsh{nonlocal} 상태 관리에서 주의할 점}

클로저를 사용하면 상태를 간단하게 캡슐화할 수 있다.

하지만 몇 가지 경계 조건을 이해해야 한다.

\begin{dhknote}
\begin{dhksmry}[centered]
1. 클로저 상태는 자동으로 불변이 되지 않는다
\end{dhksmry}
\end{dhknote}

함수 객체의 \zsh{__closure__}에 포함된 셀은 외부 바인딩에 접근하기 위한 구조다.

그 셀이 참조하는 객체가 가변 객체라면 내부 상태를 변경할 수 있다.

또한 \zsh{nonlocal}을 사용하면 바인딩 자체를 변경할 수도 있다.

따라서 클로저를 사용한다는 사실만으로 상태가 불변이거나 변경으로부터 보호된다고 볼 수는 없다.

\begin{dhknote}
\begin{dhksmry}[centered]
2. 상태 변경은 자동으로 스레드 안전하지 않다
\end{dhksmry}
\end{dhknote}

다음 연산을 생각해 보자.

\begin{dhkcode}
count += 1
\end{dhkcode}

이 연산은 기존 값을 읽고, 계산하고, 결과를 다시 바인딩하는 과정을 포함한다.

여러 실행 흐름이 동일한 클로저 상태를 동시에 변경한다면 경쟁 상태(Race Condition)가 발생할 수 있다.

따라서 클로저를 사용한다는 이유만으로 동시성 안전성이 보장되는 것은 아니다.

\begin{dhknote}
\begin{dhksmry}[centered]
3. 외부 바인딩을 유지하면 객체의 수명도 영향을 받을 수 있다
\end{dhksmry}
\end{dhknote}

클로저 셀이 큰 객체를 참조하고 있다면, 함수 객체가 유지되는 동안 그 객체도 계속 참조될 수 있다.

따라서 클로저에 불필요하게 큰 데이터를 캡처하면 메모리 사용에 영향을 줄 수 있다.

이 문제는 장기간 유지되는 콜백 함수나 데코레이터를 설계할 때 고려해야 한다.

\begin{dhknote}
\begin{dhksmry}[centered]
nonlocal은 객체의 속성 탐색과 무관하다
\end{dhksmry}
\end{dhknote}

다음 표현식을 비교해 보자.

\begin{dhkcode}
nonlocal state
state = new_state


state.value = 10
\end{dhkcode}

첫 번째는 외부 이름 \zsh{state}의 바인딩을 변경한다.

두 번째는 \zsh{state}가 참조하는 객체의 속성에 대입한다.

따라서 두 연산은 서로 다른 메커니즘이다.

객체의 속성 변경에는 일반적으로 \zsh{nonlocal}이 필요하지 않다.

\subsubsection{종합 예제 - 여러 클로저가 하나의 상태를 관리하는 구조}

이번 학습의 내용을 하나의 프로그램으로 통합해 보자.

\begin{dhkcode}
def make_account(initial_balance):
    balance = initial_balance

    def deposit(amount):
        nonlocal balance

        balance += amount

        return balance

    def withdraw(amount):
        nonlocal balance

        if amount > balance:
            raise ValueError("잔액이 부족합니다.")

        balance -= amount

        return balance

    def get_balance():
        return balance

    return deposit, withdraw, get_balance


deposit, withdraw, get_balance = make_account(1000)

print(get_balance())

print(deposit(500))

print(withdraw(300))

print(get_balance())
\end{dhkcode}

예상 출력:

\begin{dhkzsh}
1000
1500
1200
1200
\end{dhkzsh}

이 코드는 세 개의 내부 함수가 하나의 상태를 공유하는 구조다.

\begin{dhknote}
\begin{dhksmry}[centered]
객체 관계
\end{dhksmry}
\end{dhknote}

\begin{dhkdiagram}[in=6]
node n1 (1,1) :: deposit 함수 객체
node n2 (3,1) :: withdraw 함수 객체
node n3 (5,1) :: get_balance 함수 객체
node n4 (3,2) :: Cell balance
node n5 (3,3) :: int(1200)

edge n1 -> n4 :: HV
edge n2 -> n4
edge n3 -> n4 :: HV
edge n4 -> n5
\end{dhkdiagram}

여기서 \zsh{balance}는 세 함수의 자유 변수다.

\zsh{deposit}과 \zsh{withdraw}는 \zsh{nonlocal balance}를 선언했으므로 외부 바인딩을 변경할 수 있다.

반면 \zsh{get_balance}는 값을 읽기만 하므로 \zsh{nonlocal}이 필요하지 않다.

이 구조를 보면 다음 사실을 알 수 있다.

여러 함수 객체가 동일한 셀을 공유하면, 하나의 함수가 변경한 바인딩을 다른 함수도 관찰할 수 있다.

이는 클래스 인스턴스에서 여러 메서드가 동일한 인스턴스 상태를 읽고 변경하는 구조와 대응된다.

다만 이 예제는 클로저 상태 관리의 원리를 보여주기 위한 것이며, 실제 금융 거래 시스템에 필요한 동시성 제어와 영속성 등을 구현한 것은 아니다.

\subsubsection{핵심 정리와 판단 기준}

이번 학습에서 확립해야 할 핵심 개념을 정리해 보자.

\begin{dhknote}
\begin{dhksmry}[centered]
1. nonlocal은 이름의 바인딩 대상을 지정한다
\end{dhksmry}
\end{dhknote}

\begin{dhkcode}
nonlocal x
\end{dhkcode}

는 현재 함수에서 \zsh{x}를 지역 이름으로 취급하지 않고 렉시컬 외부 함수의 기존 바인딩을 사용하도록 지정한다.

\begin{dhknote}
\begin{dhksmry}[centered]
2. 셀 객체와 셀이 참조하는 객체는 구별해야 한다
\end{dhksmry}
\end{dhknote}

\begin{dhkcode}
cell is other_cell
\end{dhkcode}

은 셀 객체의 동일성을 비교한다.

반면

\begin{dhkcode}
cell.cell_contents
\end{dhkcode}

는 셀에 현재 바인딩된 객체를 조회한다.

따라서 같은 셀 객체가 시간에 따라 서로 다른 객체를 참조할 수 있다.

\begin{dhknote}
\begin{dhksmry}[centered]
3. 객체 변경과 이름 재바인딩은 다르다
\end{dhksmry}
\end{dhknote}

\begin{dhkcode}
values.append(10)
\end{dhkcode}

은 리스트 객체의 상태를 변경한다.

반면

\begin{dhkcode}
values = [10]
\end{dhkcode}

은 이름의 바인딩을 변경한다.

외부 함수의 바인딩을 변경하려면 \zsh{nonlocal}이 필요하다.

\begin{dhknote}
\begin{dhksmry}[centered]
4. 증강 대입은 이름의 스코프 판정에 영향을 준다
\end{dhksmry}
\end{dhknote}

\begin{dhkcode}
count += 1
\end{dhkcode}

은 이름에 대한 대입을 포함한다.

따라서 외부 함수의 \zsh{count}를 대상으로 하려면 \zsh{nonlocal count}가 필요하다.

가변 객체에 대한 \zsh{+=}도 이 이름 분류 규칙에서 예외가 아니다.

\begin{dhknote}
\begin{dhksmry}[centered]
5. 여러 클로저는 동일한 셀을 공유할 수 있다
\end{dhksmry}
\end{dhknote}

같은 외부 함수 호출에서 생성된 여러 내부 함수가 동일한 외부 바인딩을 참조한다면 같은 셀을 공유할 수 있다.

따라서 한 함수의 재바인딩이 다른 함수의 반환값에 영향을 줄 수 있다.

\begin{dhknote}
\begin{dhksmry}[centered]
6. 서로 다른 외부 함수 호출은 독립적인 상태를 만들 수 있다
\end{dhksmry}
\end{dhknote}

\begin{dhkcode}
counter_a = make_counter()
counter_b = make_counter()
\end{dhkcode}

두 호출에서 생성된 클로저는 서로 다른 바인딩 환경을 가질 수 있다.

따라서 독립적인 상태 관리가 가능하다.

\begin{dhknote}
\begin{dhksmry}[centered]
최종 판단 기준
\end{dhksmry}
\end{dhknote}

\begin{dhktable}
|질문|	판단 기준|
|:---|:---|
|nonlocal이 필요한가?|	외부 함수의 이름 자체를 재바인딩하는가?|
|가변 객체 변경에 nonlocal이 필요한가?|	객체 내부만 변경한다면 일반적으로 필요하지 않음|
|+=에 nonlocal이 필요한가?|	외부 이름을 대상으로 하는 증강 대입이라면 필요|
|여러 함수가 상태를 공유하는가?|	동일한 셀 또는 동일한 가변 객체를 참조하는지 확인|
|상태 변경 후 셀 객체도 바뀌는가?|	일반적인 nonlocal 재바인딩에서는 기존 셀을 유지|
|global과 nonlocal 중 무엇을 사용하는가?|	모듈 전역 바인딩인지 렉시컬 외부 함수 바인딩인지 구분|
|UnboundLocalError가 발생하는가?|	지역 이름으로 판정되었지만 읽는 시점에 바인딩되지 않았는지 확인|
|SyntaxError가 발생하는가?|	nonlocal의 대상 바인딩이 외부 함수 스코프에 존재하는지 확인|
\end{dhktable}